\providecommand{\llbracket}{[\![}
\providecommand{\rrbracket}{]\!]}
\section{A seven-rule calculus with complete progressive elimination}
\label{sec:new-calculus}

The new presentation keeps the scalar iterator and the three native formula
constructors.  It changes the meaning of an atomic closure: an atomic closure
records a progressive equation, while a compound closure retains the old
conditional-universal meaning.  A single elimination rule recovers both
one-step progress and induction from an atomic closure.  Consequently every
formula of the new language has a faithful ordinary reading; no restricted
public grammar or auxiliary-only formula class is needed.

This is a structural simplification, not a claim of an absolute minimum.  The
current Lean kernel has eight rule constructors.  The proposed kernel has
seven, and its introduction rule commutes directly with numerical
substitution.  The earlier seven-rule \texttt{WPrf} presentation also has seven
constructors, but its introduction distinguishes the literal witness $0$ from
other witnesses.  Here the distinction depends only on whether a matrix is a
literal equation, which numerical substitution cannot change.  The concrete
checker measurements, including comparisons in which this calculus is larger,
are reported separately.\cite{native,witness}

\subsection{Complete primitive specification}
\label{subsec:new-specification}

There is one numerical sort.  Terms are
\[
 t ::= x\mid 0\mid St\mid I(n,a;u.h).
\]
Only $u$ is bound in $h$; other occurrences are free parameters.  Substitution
is capture-free.  Put
\begin{equation}
 1:=S0,\qquad
 \mathsf A(a,n):=I(a,n;u.Su).
 \label{eq:new-add}
\end{equation}
Thus the first argument of $\mathsf A$ is the iteration counter and the
second is its initializer.  This argument order is essential below.

Formulas are
\begin{equation}
 F ::= (s=t)\mid Q_yF\mid D_{x,y}F.
 \label{eq:new-formulas}
\end{equation}
The binder $Q$ binds $y$, and $D$ binds the distinct variables $x,y$.
A \emph{literal equation} has outermost constructor $=$.  A \emph{compound}
formula has outermost constructor $Q$ or $D$.  Define the syntactic selector
\begin{equation}
 \eta_B(a):=
 \begin{cases}
 a,&B\text{ is a literal equation},\\
 0,&B\text{ is compound}.
 \end{cases}
 \label{eq:new-eta}
\end{equation}
Write $\Gamma\vdash_{\mathcal R}F$ for derivability.  Contexts are finite
reusable lists of formulas.  The complete primitive inventory is the following
seven rule schemas:
\[
 \frac{F\in\Gamma}{\Gamma\vdash_{\mathcal R}F}
 \quad\mathrm{Hyp},
\]
\[
 \frac{}{\Gamma\vdash_{\mathcal R}I(0,a;u.h)=a}
 \quad\mathrm{I0},\qquad
 \frac{}{\Gamma\vdash_{\mathcal R}
 I(Sn,a;u.h)=h[I(n,a;u.h)/u]}
 \quad\mathrm{IS},
\]
\begin{equation}
 \frac{\Gamma\vdash_{\mathcal R}s=t\qquad
       \Gamma\vdash_{\mathcal R}p(s)=q(s)}
      {\Gamma\vdash_{\mathcal R}p(t)=q(t)}
 \quad\mathrm{Tr}_{=},
 \label{eq:new-tr}
\end{equation}
\begin{equation}
 \frac{\Gamma\vdash_{\mathcal R}Q_yB(y)\qquad
       \Gamma\vdash_{\mathcal R}B(\eta_B(a))}
      {\Gamma\vdash_{\mathcal R}
         B(\mathsf A(\eta_B(a),n))}
 \quad\mathrm{QE}_{\mathsf A},
 \label{eq:new-qe}
\end{equation}
\begin{equation}
 \frac{\Gamma\vdash_{\mathcal R}D_{x,y}B(x,y)\qquad
       \Gamma,Q_yB(b,y)\vdash_{\mathcal R}F}
      {\Gamma\vdash_{\mathcal R}F}
 \quad\mathrm{DE},
 \label{eq:new-de}
\end{equation}
\begin{equation}
 \frac{\Gamma,B(t,\eta_B(z))\vdash_{\mathcal R}B(t,Sz)}
      {\Gamma\vdash_{\mathcal R}D_{x,y}B(x,y)}
 \quad\mathrm{DI}_{\eta}.
 \label{eq:new-di}
\end{equation}
In \eqref{eq:new-tr}, $p(v),q(v)$ are numerical term templates.  In
\eqref{eq:new-qe}, $a,n$ are arbitrary terms in the conclusion's context.
In DE, $b$ is fresh for $\Gamma,F$ and all fixed parameters of $B$.
In DI, $z$ is fresh for $\Gamma,t$ and all fixed parameters of $B$.
The displayed local assumptions are discharged.  The designated
contradiction is $0=1$.

There are no other primitive axioms, inference rules, induction schemas,
logical connectives, successor axioms, or numerical case principles.
In particular, equality transport in a compound formula is not primitive.

\begin{lemma}[Structural rules and substitution]
\label{lem:new-structural}
Weakening, exchange, contraction, renaming, assumption cut, and capture-free
numerical substitution are admissible in $\mathcal R$.
\end{lemma}
\begin{proof}
Use induction on a finite derivation, renaming eigenparameters before a
substitution or a cut.  A cut replaces each relevant hypothesis leaf by the
supplied derivation, weakened into its local context.  The only new
substitution check is
\[
 (\eta_B(a))\sigma=\eta_{B\sigma}(a\sigma).
\]
It holds because numerical substitution preserves the outermost formula
constructor.  The iterator abbreviation \eqref{eq:new-add} also commutes with
substitution.  Hence QE and DI translate to instances of the same rules;
there is no witness-literal exception to repair.
\end{proof}

\subsection{Equality and the induction bootstrap}
\label{subsec:new-bootstrap}

A useful derived rule is
\begin{equation}
 \frac{\Gamma,B(\eta_B(z))\vdash_{\mathcal R}B(Sz)}
      {\Gamma\vdash_{\mathcal R}Q_yB(y)}
 \quad\mathrm{QI}_{\eta}.
 \label{eq:new-qi}
\end{equation}
To obtain it, apply DI to $D_{x,y}B(y)$ with $x$ unused and witness $0$.
Apply DE to this result; its local hypothesis is precisely $Q_yB(y)$, which
can be exported by Hyp because it is independent of the fresh witness.

\begin{lemma}[Numerical equality and contextual equation induction]
\label{lem:new-ei}
The usual reflexivity, symmetry, transitivity, and numerical congruence rules
are derivable.  For every literal equation $E(y)$,
\begin{equation}
 \frac{\Gamma\vdash_{\mathcal R}E(0)\qquad
       \Gamma,E(z)\vdash_{\mathcal R}E(Sz)}
      {\Gamma\vdash_{\mathcal R}E(t)}
 \quad(z\text{ fresh})
 \label{eq:new-ei}
\end{equation}
is admissible with arbitrary retained side assumptions $\Gamma$.
\end{lemma}
\begin{proof}
Let $J(a)=I(0,a;u.0)$.  I0 gives $J(a)=a$; transport this equation along
itself with template $v=a$ to obtain $a=a$.  The usual equational templates
then give symmetry, transitivity, and numerical congruence.

For induction, the step premise and \eqref{eq:new-qi} give $Q_yE(y)$.
Use \eqref{eq:new-qe} with source $a=0$ and initializer $n=t$.  Its conclusion
is $E(I(0,t;u.Su))$.  I0 and purely equational transport give $E(t)$.
No identity $I(t,0;u.Su)=t$ has been assumed.  Using the counter $t$ and
initializer $0$ in the elimination rule instead would not justify this
bootstrap.
\end{proof}

\begin{lemma}[Elementary iterator identities]
\label{lem:new-addition}
The new calculus proves, with arbitrary parameters,
\begin{equation}
 \begin{aligned}
 \mathsf A(0,n)&=n,&
 \mathsf A(Sa,n)&=S\mathsf A(a,n),\\
 \mathsf A(a,0)&=a,&
 \mathsf A(a,Sn)&=S\mathsf A(a,n),\\
 \mathsf A(a,n)&=\mathsf A(n,a),&
 \mathsf A(a,1)&=Sa.
 \end{aligned}
 \label{eq:new-add-identities}
\end{equation}
\end{lemma}
\begin{proof}
The first row consists of computation instances.  Induct on $a$ for
$\mathsf A(a,0)=a$.  Induct on $a$ again for the initializer-successor
identity: its base is $Sn=Sn$, and its step follows by applying $S$ to the
induction hypothesis.  For commutativity, induct on $a$.  The base is
$n=\mathsf A(n,0)$; the step is
\[
 \mathsf A(Sa,n)=S\mathsf A(a,n)
               =S\mathsf A(n,a)=\mathsf A(n,Sa).
\]
Finally $\mathsf A(a,1)=S\mathsf A(a,0)=Sa$.
All inductions are instances of \eqref{eq:new-ei} on literal equations.
\end{proof}

An atomic closure therefore has the derived one-step elimination
\begin{equation}
 Q_yE(y),\ E(a)\ \vdash_{\mathcal R}\ E(Sa).
 \label{eq:new-step}
\end{equation}
Take initializer $1$ in QE and use \eqref{eq:new-add-identities}.
For a compound matrix, QE at source $0$ and initializer $t$ has the ordinary
conditional-closure effect $Q_yB,B(0)\vdash B(t)$, once transport through
$B$ has been derived.

\begin{lemma}[Full transport and homogeneous binder congruence]
\label{lem:new-transport}
For every formula template $F(v)$,
\begin{equation}
 s=t,\ F(s)\ \vdash_{\mathcal R}\ F(t).
 \label{eq:new-full-tr}
\end{equation}
Moreover, if $B$ and $C$ are uniformly interderivable and have the same root
class (both equations or both compound), their $Q$ closures are
interderivable, and their $D$ closures are interderivable.
\end{lemma}
\begin{proof}
For transport, induct on $F$.  Equations use the primitive rule.  For
$Q_yE(v,y)$ with atomic $E$, introduce the target closure by QI.
Reverse-transport its local $E(t,z)$ to $E(s,z)$, use
\eqref{eq:new-step} for the source closure, and transport the resulting
$E(s,Sz)$ to $E(t,Sz)$.

For $Q_yB(v,y)$ with compound $B$, introduce the target closure from
$B(t,0)$.  The induction hypothesis converts this to $B(s,0)$.
QE for the source closure at initializer $Sz$ gives
$B(s,I(0,Sz;u.Su))$.  The induction hypothesis applied inside the strictly
smaller matrix $B$ removes the I0 redex and changes the parameter $s$ to $t$.
This proves the QI target $B(t,Sz)$.

For $D_{x,y}B(v,x,y)$, eliminate the source binder at a fresh witness $b$.
Introduce the target at that witness.  Convert its local assumption
$B(t,b,\eta_B(z))$ to $B(s,b,\eta_B(z))$ by the induction hypothesis.
At an atomic matrix use \eqref{eq:new-step}; at a compound matrix use QE
with source $0$, initializer $Sz$, and I0.  Transport the successor result
back to the parameter $t$.  Close DI and DE.  Every invoked formula
transport concerns the strictly smaller matrix, so this is not circular.

For homogeneous congruence, use the same constructions, replacing numerical
transport by the two given conversions between $B$ and $C$.  Their root
classes agree, so their local DI assumptions use the same value of $\eta$.
For $D$, eliminate the source at a fresh witness and introduce the target
at that witness.  The argument works uniformly under arbitrary additional
retained assumptions.
\end{proof}

The root-class qualification is substantive.  The new $Q$ is not claimed to
preserve arbitrary equivalence when one matrix is an equation and the other
is compound.  A progressive assertion and a conditional-universal assertion
are different assertions even when their matrices are pointwise equivalent.
The translations below respect this boundary explicitly.

\subsection{Ordinary soundness}
\label{subsec:new-soundness}

Use the pinned repository's scalar theory $T_{\rm eq}$: ordinary
intuitionistic first-order logic with actual equality, the two iterator
computations, numerical coverage
\[
 \forall x\,(x=0\lor\exists y\,x=Sy),
\]
separation $(0=1)\to\bot$, and contextual induction for literal equations.
The repository proves its whole-domain bi-interpretability with iPRA.\cite{bridge}

Define the new ordinary reading recursively by
\begin{equation}
 \begin{aligned}
 \llbracket s=t\rrbracket&=(s=t),\\
 \llbracket Q_yE(y)\rrbracket
   &=\forall y\,(E(y)\to E(Sy))
       &&(E\text{ a literal equation}),\\
 \llbracket Q_yB(y)\rrbracket
   &=\llbracket B(0)\rrbracket\to
                  \forall y\,\llbracket B(y)\rrbracket
       &&(B\text{ compound}),\\
 \llbracket D_{x,y}B\rrbracket
   &=\exists x\,\llbracket Q_yB\rrbracket.
 \end{aligned}
 \label{eq:new-reading}
\end{equation}
This translation preserves all free numerical variables and commutes with
numerical substitution.  It is a finite syntactic translation, not a truth
predicate on numerical formula codes.

\begin{theorem}[Soundness]
\label{thm:new-soundness}
If $\Gamma\vdash_{\mathcal R}F$, then
$\llbracket\Gamma\rrbracket\vdash_{T_{\rm eq}}\llbracket F\rrbracket$.
\end{theorem}
\begin{proof}
Hyp, computations, equational transport, and DE have their ordinary meanings.
Atomic DI is universal introduction of progress, followed by existential
introduction at its witness.  Compound DI is conditional universal
introduction using numerical coverage, followed by existential introduction.

For atomic QE, retain
$H=\forall y(E(y)\to E(Sy))$ and $E(a)$.
Equation induction on $n$ gives $E(\mathsf A(n,a))$: the base is $E(a)$,
and the step uses $H$ at $\mathsf A(n,a)$ and IS.  The pure equational
proofs in Lemma~\ref{lem:new-addition} also hold in $T_{\rm eq}$, so
commutativity gives $E(\mathsf A(a,n))$.  For compound QE the minor is
$\llbracket B(0)\rrbracket$; its conditional-universal major directly
supplies the target instance.  Thus every rule translates soundly.
\end{proof}

\subsection{A precise comparison with the existing native calculus}
\label{subsec:new-native-comparison}

Let $\mathcal N$ denote exactly the current native Eq/$Q$/$D$ calculus
\texttt{NPrf}.  Its $Q_yB$ always means $B(0)\to\forall yB(y)$; its direct
progress rule introduces $Q_yE$ only for equations.  Its terms and numerical
equality are identical to those of $\mathcal R$.

For $E\equiv(s=t)$ define
\begin{equation}
 \begin{aligned}
 C(v,a,b)&:=I(v,a;u.b),\\
 G_vE&:=\bigl(C(v,0,s)=C(v,0,t)\bigr),\\
 L(E)&:=D_{r,v}G_vE,\\
 P(F)&:=\begin{cases}L(F),&F\text{ an equation},\\F,&F\text{ compound}.
              \end{cases}
 \end{aligned}
 \label{eq:new-padding}
\end{equation}
All newly displayed binders are fresh; the update in $C$ ignores $u$, and
$r$ is unused in $L(E)$.

\begin{lemma}[Padding]
\label{lem:new-padding}
In $\mathcal R$, $L(E)$ and $E$ are interderivable under arbitrary contexts.
Consequently $P(F)$ is compound and interderivable with $F$.
\end{lemma}
\begin{proof}
Computations and equality prove $G_0E$ and
$G_{Sz}E\dashv\vdash E$.  From $E$, introduce $L(E)$ by atomic DI, ignoring
its local induction assumption and using $E$ for the successor target.
Conversely, DE on $L(E)$ gives $Q_vG_vE$.  Supply the theorem $G_0E$ and
apply QE with source $0$, initializer $1$.  I0 gives $G_1E$, which reduces
to $E$.  This proof uses no ordinary logical macro.
\end{proof}

Define a translation $T:\mathcal N\to\mathcal R$ by
\begin{equation}
 \begin{aligned}
 T(s=t)&=(s=t),\\
 T(Q_yB)&=Q_yP(TB),\\
 T(D_{x,y}B)&=D_{x,y}P(TB).
 \end{aligned}
 \label{eq:new-T}
\end{equation}
It fixes terms, commutes with numerical substitution, and places a compound
matrix immediately below every translated binder.

\begin{theorem}[Translation into the new calculus]
\label{thm:new-T}
Every $\mathcal N$ derivation translates effectively:
\[
 \Gamma\vdash_{\mathcal N}A
 \quad\Longrightarrow\quad
 T\Gamma\vdash_{\mathcal R}TA.
\]
\end{theorem}
\begin{proof}
Induct on the old derivation.  Hyp, computations, and equational transport
translate literally.  Old QE becomes compound new QE: pad the minor,
specialize at the old target using source $0$, remove the I0 redex, and
unpad.  In DE the opened new assumption is exactly
$Q_yP(TB)(b,y)$, the translation of the old opened $Q$ assumption.
For case DI, unpad the local base, apply the translated premise, and pad
the successor conclusion; the new matrix is compound.

For old direct Q-progress at $E$, the translated target is $Q_yL(E(y))$.
Introduce it by compound QI.  Its local base $L(E(0))$ supplies $E(0)$.
Apply the recovered contextual equation induction to the translated step,
then pad $E(Sz)$ to $L(E(Sz))$.  This discharges the only nonliteral rule
translation.  Each construction is finite and preserves eigenvariable
conditions after fresh renaming.
\end{proof}

\subsection{Logical macros and the order of the completeness proof}
\label{subsec:new-image-logic}

The old ordinary logical macros will be denoted
$\bot_{\mathcal N}$, $\to_{\mathcal N}$, $\wedge_{\mathcal N}$,
$\vee_{\mathcal N}$, $\exists_{\mathcal N}$, and $\forall_{\mathcal N}$.
For specificity they are the macros of the pinned
\texttt{NativeLogic.lean}:\cite{native}
\begin{equation}
 \begin{aligned}
 \bot_{\mathcal N}&=(0=1),\\
 A\to_{\mathcal N}B&=Q_y\operatorname{sel}(y,A,B),\\
 A\wedge_{\mathcal N}B
   &=Q_y\operatorname{gate}
           (y,\operatorname{sel}(\operatorname{par}(y),A,B)),\\
 \exists_{\mathcal N}x\,A(x)&=D_{x,y}\operatorname{gate}(y,A(x)),\\
 A\vee_{\mathcal N}B&=\exists_{\mathcal N}v\,
                                   \operatorname{sel}(v,A,B),\\
 \forall_{\mathcal N}x\,A(x)&=A(0)\wedge_{\mathcal N}Q_xA(x).
 \end{aligned}
 \label{eq:new-old-macros}
\end{equation}
Here $\operatorname{par}(y)=I(y,0;u.C(u,1,0))$.
The old gate acts on equations as $G_v$ in \eqref{eq:new-padding}.
The exact finite recursive definitions of the gate and selector, including
their behavior under $Q$ and $D$, are stated in Appendix~\ref{app:old-macros}.
The selector has zero endpoint $A$ and successor endpoint $B$.
All these operations concern the \emph{old} uniform conditional calculus.
The pinned repository supplies their capture-free definitions and derived
rules.

\begin{lemma}[Logic on image formulas, with arbitrary side context]
\label{lem:new-image-logic}
For old formulas $A,B$, the formulas
$T(A\to_{\mathcal N}B)$, $T(A\wedge_{\mathcal N}B)$, and the analogous
translated macros have their ordinary natural-deduction rules on principal
formulas $TA,TB$.  These derived rules permit an arbitrary retained
$\mathcal R$ context and arbitrary supplied $\mathcal R$ premise proofs.
Image existential elimination also permits an arbitrary new conclusion,
not necessarily an image formula.
\end{lemma}
\begin{proof}
Translate the finite \emph{derived-rule constructions}, rather than merely
closed old theorems.  Their nonlogical steps are replaced by the local
simulations in Theorem~\ref{thm:new-T}.  Their open premise leaves are filled
with the supplied new proofs; assumption cut and weakening
(Lemma~\ref{lem:new-structural}) retain any additional new context.  The
construction never examines the syntax of passive side assumptions.
Thus this is a uniform rule simulation, not an inference from
closed-theorem preservation alone.

For the important existential exception, the old elimination construction
opens its defining $D$ by DE, obtains the witness instance from the opened
gated closure, and cuts this instance into the supplied minor derivation.
It neither analyzes nor introduces the minor's conclusion.  The translated
construction consequently works for every new conclusion $F$, using the
same freshness condition.  This is precisely the structure of the pinned
\texttt{NPrf.exE} proof.
\end{proof}

We do \emph{not} yet assert that the old selectors or logical macros can be
applied unchanged to arbitrary mixed formulas.  Such an assertion would
incorrectly ignore the root distinction in $Q$.  The next normalization
argument first proves that every mixed formula is equivalent to an image.

\subsection{Both formula composites for the entire language}
\label{subsec:new-composites}

Define $J:\mathcal R\to\mathcal N$, fixing all terms and equations, by
\begin{equation}
 \begin{aligned}
 J(Q_yE)&=\forall_{\mathcal N}y\,
                (E(y)\to_{\mathcal N}E(Sy)),
                    &&E\text{ an equation},\\
 J(Q_yB)&=Q_yJB,     &&B\text{ compound},\\
 J(D_{x,y}E)&=\exists_{\mathcal N}x\,
             \forall_{\mathcal N}y\,
                (E(x,y)\to_{\mathcal N}E(x,Sy)),
                    &&E\text{ an equation},\\
 J(D_{x,y}B)&=D_{x,y}JB,&&B\text{ compound}.
 \end{aligned}
 \label{eq:new-J}
\end{equation}
This is a fixed finite syntax transformation.  It maps compound formulas to
compound formulas, preserves free variables, and commutes with numerical
substitution.  Its ordinary reading in $T_{\rm eq}$ is provably equivalent to
\eqref{eq:new-reading}, by the established meanings of the old logical macros.

\begin{lemma}[Atomic closure normalization]
\label{lem:new-atomic-normalization}
For every equation $E(x,y)$, under arbitrary retained contexts,
\begin{equation}
 \begin{aligned}
 Q_yE(x,y)&\dashv\vdash_{\mathcal R}
       T\bigl(\forall_{\mathcal N}y
           (E(x,y)\to_{\mathcal N}E(x,Sy))\bigr),\\
 D_{x,y}E(x,y)&\dashv\vdash_{\mathcal R}
       T\bigl(\exists_{\mathcal N}x\,
             \forall_{\mathcal N}y
           (E(x,y)\to_{\mathcal N}E(x,Sy))\bigr).
 \end{aligned}
 \label{eq:new-atomic-normalization}
\end{equation}
\end{lemma}
\begin{proof}
Write $H(x)$ for the image formula on the first line's right.
From $Q_yE(x,y)$ and the temporary assumption $E(x,z)$, derived one-step
elimination gives $E(x,Sz)$.  The image implication and universal
introductions of Lemma~\ref{lem:new-image-logic} give $H(x)$.  Their
principal component formulas are just equations, which are already fixed
by $T$; no full mixed-language logical rule has been assumed.
Conversely, under $H(x)$ and $E(x,z)$, image universal specialization and
implication elimination give $E(x,Sz)$; atomic QI gives $Q_yE(x,y)$.

For $D$, eliminate the left-hand binder at fresh $b$, normalize its local
$Q_yE(b,y)$ to $H(b)$ by the first part, introduce the image existential
at $b$, and close DE.  In the reverse direction, image existential
elimination supplies a fresh $b$ and $H(b)$.  Under $E(b,z)$ its image
universal and implication rules give $E(b,Sz)$, so atomic DI at witness
$b$ gives $D_{x,y}E$.  Close image existential elimination.  Its conclusion
is allowed to be this arbitrary new formula by the last assertion of
Lemma~\ref{lem:new-image-logic}.
\end{proof}

\begin{theorem}[Exact native-language round trips]
\label{thm:new-roundtrips}
For every old formula $A$ and every new formula $F$,
\begin{equation}
 A\dashv\vdash_{\mathcal N}JT(A),\qquad
 F\dashv\vdash_{\mathcal R}TJ(F).
 \label{eq:new-roundtrips}
\end{equation}
These equivalences are uniform under arbitrary additional retained contexts.
\end{theorem}
\begin{proof}
For the old-side composite, note first that $JL(E)$ is old-equivalent to
$E$.  It asserts existence of an unused witness and progressivity of the
guarded equation $G_vE$.  From $E$, every positive guard instance follows,
so the progress implications hold.  Conversely $G_0E$ is a theorem and
progress at $0$ yields $G_1E$, hence $E$.  This uses ordinary logical
macros in $\mathcal N$, whose correctness is already established.
Now induct on the old formula.  Every binder created by $T$ has a compound
matrix, so $J$ gives its old conditional interpretation.  The preceding
padding equivalence and the induction hypothesis apply under the old
$Q$ and $D$ congruence rules.

For the new-side composite, induct on $F$.  Equations are unchanged.
Atomic $Q$ and $D$ cases are exactly
Lemma~\ref{lem:new-atomic-normalization}.  Suppose the matrix $B$ is
compound.  The induction hypothesis gives $B\dashv\vdash TJB$.
Both $B$ and $TJB$ are compound, since $J$ and $T$ preserve compound roots.
Homogeneous congruence (Lemma~\ref{lem:new-transport}) therefore gives
\[
 Q B\dashv\vdash Q(TJB),\qquad
 D B\dashv\vdash D(TJB).
\]
Because $TJB$ is compound, no padding is inserted on these right sides;
they are exactly the required $TJ$ composites.  There is no appeal to
congruence across unlike root classes.  The induction proves normalization
for every formula before ordinary logic is extended to arbitrary new
formulas, avoiding a circular completeness argument.
\end{proof}

\paragraph{All ordinary logical operations in the new language.}
Now, and only now, define for arbitrary new formulas
\begin{equation}
 \begin{aligned}
 \bot_{\mathcal R}&=(0=1),\\
 F\mathbin{\diamond_{\mathcal R}}G
   &:=T\bigl(JF\mathbin{\diamond_{\mathcal N}}JG\bigr)
       &&(\diamond\in\{\to,\wedge,\vee\}),\\
 \exists_{\mathcal R}x\,F&:=T(\exists_{\mathcal N}x\,JF),\\
 \forall_{\mathcal R}x\,F&:=T(\forall_{\mathcal N}x\,JF).
 \end{aligned}
 \label{eq:new-full-macros}
\end{equation}
The round trips convert arbitrary new principal formulas to their images;
Lemma~\ref{lem:new-image-logic} then supplies the ordinary rules, and the
round trips convert their conclusions back.  For disjunction or existential
elimination with a new conclusion, normalize that conclusion first if
necessary.  All eigenvariable conditions are preserved because $J,T$
preserve free numerical variables.

This gives intuitionistic natural deduction for the entire new language,
including explosion, ordinary numerical coverage, and the quantified
rules.  It adds no primitive logical axioms.  In particular, the system does
not use excluded middle for arbitrary quantified formulas.

\subsection{Bi-interpretability, consistency, and the inherited arithmetic}
\label{subsec:new-final-bridge}

\begin{theorem}[Full consequence equivalence and bi-interpretability]
\label{thm:new-biinterpretability}
The translations $T,J$ preserve derivations in both directions and have the
round trips \eqref{eq:new-roundtrips}.  The ordinary realization of
$\mathcal R$ is equiconsistent and bi-interpretable with full first-order
iPRA.  Both numerical domains are whole, numerical equality is actual
equality, and the numerical composite isomorphisms have graph $x=y$.
\end{theorem}
\begin{proof}
Preservation by $T$ is Theorem~\ref{thm:new-T}.  For $J$, translate a new
derivation into $T_{\rm eq}$ using Theorem~\ref{thm:new-soundness}; its
formulas are provably equivalent to the ordinary readings of their $J$
translations.  The pinned
native completeness and native formula round-trip theorems translate this
ordinary derivation back into $\mathcal N$ and identify each resulting
formula with its $J$ translation.  Equivalently, one may translate each new
rule directly using the old ordinary macros and equation induction; the
only added arithmetic calculation is the QE soundness proof above.
The two exact formula composites are Theorem~\ref{thm:new-roundtrips}.

The remaining scalar-to-iPRA arithmetic is inherited, not reproved by the
new elimination rule.  The repository's chain constructs a scalar
predecessor, a term reflecting actual equality, quantifier-free induction,
pairing, and a compiler for each fixed primitive-recursive description.
The reverse interpretation expresses a fixed scalar iterator by primitive
recursion.  Equation induction proves the two numerical compilation round
trips.  These are the contents of the pinned \texttt{biInterp},
\texttt{native\_iff\_pra}, \texttt{pra\_iff\_native},
\texttt{native\_roundTrip}, and \texttt{pra\_native\_roundTrip} results.\cite{native,bridge}

Composing that existing bi-interpretation with $T,J$ proves the assertion.
Our two translations leave every scalar term and the numerical domain
unchanged, so they add neither a cut nor a quotient to the existing
identity graphs.  The designated contradiction $0=1$ is fixed.  Effective
finite proof translations therefore give equiconsistency.  This is an
external equivalence theorem, not an internal assertion of iPRA's own
consistency.
\end{proof}

\paragraph{What changes compared with the earlier progressive proposal.}
The repository already contains a presentation with a progressive binder
on every matrix, separate one-step and equation-iteration eliminations,
counter-aware primitive recursion, and an explicit Boolean-cases rule.\cite{progressive}
The progressive reading itself is therefore inherited motivation.  The
present contribution is the root-dependent complete reading with the scalar
iterator unchanged, the single reversed-orbit specialization
\eqref{eq:new-qe}, the removal of a separate case principle, and both
round trips for the entire Eq/$Q$/$D$ formula language.  The mixed reading
without the stronger elimination would lose information at atomic closures;
its soundness alone would not establish the theorem above.

\paragraph{Verification status.}
The new metatheorems in this section are written proofs.  They are not claimed
to have been checked in Lean or Rocq.  The existing Lean theorem names identify
the precise inherited bridge.  Executable certificates and checker tests
support concrete implementations; they do not replace the universal proof
or establish a global lower bound on presentation size.
